
ce noyau est contenu dans celui de l'homomorphisme composé :
\[
\mathbb Z[T_1,X_1,\ldots,X_d]\longrightarrow
\bigoplus_{n\ge0}J^n\longrightarrow A,
\]
le dernier étant défini en degré \(n\) par l'injection de \(J^n\) dans \(A\). Or le noyau de l'homomorphisme composé est l'idéal principal engendré par \((X_1-T_1)\), qui pour des raisons d'intégrité ne contient pas d'éléments homogènes non nuls. Donc le premier homomorphisme est injectif, ce qui achève la démonstration dans le cas considéré.

Revenons au cas général où \(A\) est quelconque ; on peut encore supposer \(\underline{E}\) libre sur \(X\). Affectons l'indice \(0\) au cas particulier qui vient d'être traité ; alors d'après la partie i) de la proposition, la situation considérée est obtenue par le changement de base \(A_0\to A\) à partir de la situation étudiée précédemment, l'homomorphisme \(A_0\to A\) étant défini en prenant une base de \(E\), son image \(f_1,\ldots,f_d\) dans \(J\), et en envoyant \(T_i\) sur \(f_i\). Posons \(B_0=\left(\bigoplus_{n\ge0}J_0^n\right)_{(T_1^{(1)})}\), \(B=\left(\bigoplus_{n\ge0}J^n\right)_{(f_1^{(1)})}\) ; alors \(B=B_0\otimes_{A_0}A\) ; le complexe de Koszul défini par \(v:\underline{F}\to\underline{O}_P\) se déduit par changement de base du complexe de Koszul défini par \(v_0\) ; ce dernier étant une résolution projective de \(B_0\) (en se plaçant sur un ouvert convenable) les groupes d'homologie de \(K_\bullet(v)\) sont sur cet ouvert les \(\Tor_i^{A_0}(A,B_0)\). Ces derniers sont les composantes de degré \(0\) des \(\Tor_i^{A_0}\left(A,\left(\bigoplus_{n\ge0}J_0^n\right)_{T_1^{(1)}}\right)=\left(\bigoplus_{n\ge0}\Tor_i^{A_0}(A,J_0^n)\right)_{T_1^{(1)}}\) ; or les \(\Tor_i^{A_0}(A,J_0^n)\) sont nuls pour \(i\ge1\) d'après les relations 1.2, donc \(K_\bullet(v)\) est une résolution de \(B\), ce qui montre que \(j\) est régulière.

iii) La propriété étant locale sur \(X\), on peut encore supposer \(X\) affine et \(\underline{E}\) libre de type fini. On reprend alors les notations de ii).

On considère d'abord le cas standard de ii). L'anneau \(A_0\) est alors noethérien et de dimension cohomologique finie, et l'algèbre \(\underline{\underline{S}}(E_0)\)
égale%
